\title{Search Smooths Discontinuities in\\ Discrete/Continuous Problems}

\author{
Christopher D.\ Carroll\thanks{Carroll: Department of Economics, Johns Hopkins University,
\url{http://econ.jhu.edu/people/ccarroll/}, \texttt{ccarroll@jhu.edu}.}\\
JHU
\and
Fedor Iskhakov\thanks{Iskhakov: Stony Brook University.}\\
Stony Brook University
\and
Shujaat Khan\thanks{Khan: IMF.}\\
IMF
\and
Adrian Monninger\thanks{Monninger: Department of Economics, Texas A\&M University, \url{https://sites.google.com/view/adrianmonninger}, \texttt{amonninger@tamu.edu}.}\\
Texas A\&M
}

\date{\today}

\maketitle

\begin{abstract}
When used to solve models with both discrete and continuous choices, the Endogenous Gridpoint Method (EGM), in addition to finding the optimal solution, can also generate points in the solution that are suboptimal. This paper shows that incorporating search into such problems can smooth out the discontinuities in the decision rules, which are responsible for producing the suboptimal points. With ample search frictions, the suboptimal points can be eliminated entirely. In addition to providing greater control over the smoothness of the problem, search offers the added benefit of rationalizing the smoothness into the agent's choice problem.

\medskip
\noindent\textbf{Keywords:} Consumption, Computational Economics, Discrete Continuous Choice, Endogenous Gridpoint Method, Search

\noindent\textbf{JEL codes:} C13, C63

\medskip
\noindent\textbf{Repo:} \url{https://github.com/AMonninger/DCEGM_Search}\\
\noindent\textbf{PDF:} \url{https://github.com/AMonninger/DCEGM_Search/Text/v3/dcegmSearch_v3.pdf}
\end{abstract}

\bigskip
\noindent\footnotesize The computational results in this paper were constructed using tools in the Econ-ARK/HARK toolkit. The toolkit can be cited by its digital object identifier 10.5281/zenodo.1001068, as is done in the paper's own references as \citet{Carroll2017-hk}.
Thanks to Matthew White for initial discussions of this subject.
\normalsize
